\section{Results}

We evaluate three models under three memory policies on 99 SWE-Bench Verified
instances with a 64K context budget. Within a model every arm shares the same
sampling configuration and budget, so the only variable is the memory policy.
A1 is plain ReAct with no memory. A2 compresses on a harness threshold. A3 is
ACM's untrained Base agent, which is given the context management tools and
chooses when to use them.

\paragraph{Memory improves coverage rather than accuracy.}
Pass@1 rises for every model (Table~\ref{tab:memory-arms}). The gain comes
almost entirely from instances the baseline lost. Splitting each model's
instances by whether A1 ran out of context, A1 resolves none of the instances it
lost, while A2 and A3 resolve 11\% to 27\% of the same instances. On the
instances A1 already handled the arms are within six points of each other.
Memory therefore turns impossible instances into merely hard ones, and it does
not damage runs that never needed it.

\paragraph{Compression predicts failure without causing it.}
Instances that compressed resolve at 3\% to 28\%. Instances that did not
compress resolve at 37\% to 63\%. This is selection rather than harm. An
instance reaches the compression trigger because it is long, and long instances
are harder. The correct counterfactual is the baseline, which scores zero on
exactly those instances.

\paragraph{No model manages its own context.}
All 216 compressions were either triggered by the harness or forced at the 95\%
threshold. None were voluntary, and \texttt{query\_memory} was invoked once
across all nine runs. This holds with both tools registered as real
function-calling tools, so it is not an artifact of how they are exposed. An
earlier configuration that asked models to issue \texttt{manage\_context} as a
bash string was ignored entirely, which understated forced compliance.

\paragraph{The bottleneck moves from context to turns.}
Memory removes out-of-context failures and replaces them with turn-limit
failures. For Qwen3.5-9B, out-of-context failures fall from 66 to 5 while
turn-limit failures rise from 3 to 45. MiMo loses 69 of 99 instances to the
250-turn limit under A2. ACM imposes no turn cap, so this limit is ours, and it
truncates the memory arms more often than the baseline because memory exists to
extend the run.

\paragraph{Comparison with ACM.}
Our absolute numbers are lower than theirs, 0.172 against 0.489 for ReAct on the
same base model. We expose a single bash tool against their three-tool scaffold,
impose a 64K budget against their 128K and evaluate on a subset selected for
context pressure. The effect size moves the other way. Their ReAct to Base gap
is 3.9\%, while ours ranges from 31\% to 118\%. At 128K, SWE-Bench rarely fills
the window, so the mechanism has little to do. How much context management is
worth appears to depend on how tight the budget is, which a single operating
point cannot show.
